% ============================================================================
%  INFORME TÉCNICO — PLAN DE PRUEBAS UNITARIAS
%  Sistema AliGest — Gestión de Expensas Residenciales
%  Universidad de las Fuerzas Armadas ESPE — Ingeniería en Software
%  Asignatura: Análisis y Diseño de Software
%  Compilar en Overleaf con: pdfLaTeX
% ============================================================================
\documentclass[12pt,a4paper]{article}

% ---------- Codificación e idioma ----------
\usepackage[utf8]{inputenc}
\usepackage[T1]{fontenc}
\usepackage[spanish,es-tabla]{babel}

% ---------- Página y tipografía ----------
\usepackage[top=2.5cm,bottom=2.5cm,left=3cm,right=2.5cm]{geometry}
\usepackage{setspace}
\usepackage{lmodern}
\usepackage{microtype}

% ---------- Gráficos y tablas ----------
\usepackage{graphicx}
\usepackage{booktabs}
\usepackage{array}
\usepackage{tabularx}
\usepackage{longtable}
\usepackage{float}
\usepackage{caption}
\captionsetup{font=small,labelfont=bf}

% ---------- Listas, colores, enlaces ----------
\usepackage{enumitem}
\usepackage{xcolor}
\definecolor{espeverde}{RGB}{0,102,51}
\usepackage[colorlinks=true,linkcolor=black,citecolor=black,urlcolor=espeverde]{hyperref}

% ---------- Código fuente (bloques de implementación) ----------
% NOTA: paquete agregado sobre la plantilla base para poder incluir los
% fragmentos de código JavaScript/Jest citados a lo largo del informe.
\usepackage{listings}
\lstdefinestyle{codigo}{
    basicstyle=\ttfamily\footnotesize,
    breaklines=true,
    breakatwhitespace=false,
    backgroundcolor=\color{gray!8},
    frame=single,
    rulecolor=\color{gray!45},
    showstringspaces=false,
    tabsize=2,
    columns=fullflexible,
    keywordstyle=\color{blue!60!black}\bfseries,
    commentstyle=\color{gray!55!black}\itshape,
    stringstyle=\color{espeverde!70!black},
    numbers=left,
    numberstyle=\tiny\color{gray!70},
    numbersep=6pt,
    xleftmargin=12pt
}
\lstset{style=codigo, language=JavaScript}

% ---------- Encabezado y pie de página ----------
\usepackage{fancyhdr}
\pagestyle{fancy}
\fancyhf{}
\fancyhead[L]{\small Plan de Pruebas Unitarias — AliGest}
\fancyhead[R]{\small ESPE — Análisis y Diseño}
\fancyfoot[C]{\thepage}
\renewcommand{\headrulewidth}{0.4pt}
\renewcommand{\footrulewidth}{0pt}

% ---------- Interlineado ----------
\onehalfspacing

% Columna centrada de ancho fijo para tablas
\newcolumntype{C}[1]{>{\centering\arraybackslash}p{#1}}

\begin{document}

% ============================================================================
%  PORTADA
% ============================================================================
\begin{titlepage}
\centering
\vspace*{0.5cm}

\includegraphics[width=0.28\textwidth]{logo_espe.png}\\[0.4cm]
% NOTA (Figura de portada): reemplace "logo_espe.png" por el escudo oficial
% de la Universidad de las Fuerzas Armadas ESPE (formato PNG con fondo
% transparente). Si no dispone del logo, puede comentar esta línea con "%".

{\LARGE \textbf{UNIVERSIDAD DE LAS FUERZAS ARMADAS}}\\[0.15cm]
{\LARGE \textbf{ESPE}}\\[0.8cm]

{\large \textbf{Ingeniería en Software}}\\[0.6cm]

\rule{0.85\textwidth}{0.5pt}\\[0.6cm]
{\Large \textbf{INFORME TÉCNICO}}\\[0.35cm]
{\huge \textbf{Plan de Pruebas Unitarias}}\\[0.2cm]
{\Large Sistema AliGest — Gestión de Expensas Residenciales}\\[0.35cm]
\rule{0.85\textwidth}{0.5pt}\\[1.4cm]

\begin{flushleft}
\hspace{2cm}\textbf{Asignatura:} \hspace{0.6cm} Análisis y Diseño\\[0.35cm]
\hspace{2cm}\textbf{NRC:} \hspace{2.05cm} 30723\\[0.35cm]
\hspace{2cm}\textbf{Docente:} \hspace{1.35cm} Ing. Jenny Ruiz\\[0.35cm]
\hspace{2cm}\textbf{Integrantes:} \hspace{0.75cm} Santi Jeancarlo\\[0.15cm]
\hspace{2cm}\phantom{\textbf{Integrantes:}} \hspace{0.75cm} Erick Tufiño\\[0.15cm]
\hspace{2cm}\phantom{\textbf{Integrantes:}} \hspace{0.75cm} Andrés Cedeño\\[0.35cm]
\hspace{2cm}\textbf{Fecha:} \hspace{1.7cm} 09 de julio de 2026\\
\end{flushleft}

\vfill
{\large \textbf{Sangolquí — Ecuador}}\\[0.2cm]
{\large \textbf{2026}}

\end{titlepage}

% ============================================================================
%  ÍNDICES
% ============================================================================
\pagenumbering{roman}
\setcounter{page}{2}

\renewcommand{\contentsname}{Índice de Contenidos}
\tableofcontents
\clearpage

\listoftables
\clearpage

% ============================================================================
%  CUERPO DEL INFORME
% ============================================================================
\pagenumbering{arabic}
\setcounter{page}{1}

\section{Objetivo de las Pruebas}

\subsection{Objetivo General}

Validar que los requisitos funcionales del sistema \textbf{AliGest} —plataforma de
gestión de expensas residenciales del conjunto habitacional \emph{``La Primavera''}—
cumplan con las especificaciones establecidas en el backlog del proyecto,
garantizando el correcto funcionamiento de:

\begin{itemize}[nosep]
    \item Autenticación, registro y recuperación de acceso de usuarios.
    \item Configuración de perfiles y control de permisos por rol (Administrador / Copropietario).
    \item Administración de copropietarios: importación masiva desde Excel, consulta, modificación y eliminación controlada.
    \item Registro, validación y auditoría de pagos de expensas.
    \item Cálculo automático del recargo por mora sobre deudas vencidas.
    \item Generación de comprobantes digitales inmutables y notificaciones automáticas al residente.
\end{itemize}

\subsection{Objetivos Específicos}

\begin{enumerate}[nosep]
    \item \textbf{Validar objetos de valor y modelo de dominio:} verificar que \texttt{Cedula},
    \texttt{Telefono} y \texttt{Usuario} apliquen correctamente sus reglas de negocio e
    invariantes (algoritmo Módulo 10, formato internacional ecuatoriano, políticas de
    contraseña).
    \item \textbf{Validar servicios de negocio:} asegurar que \texttt{AuthService},
    \texttt{CopropietarioService} y \texttt{PagoService} gestionen correctamente la lógica
    de negocio, las validaciones cruzadas contra la base de datos y el manejo de errores.
    \item \textbf{Validar los patrones de diseño aplicados:} confirmar que el patrón
    \emph{Strategy} (cálculo de mora), el patrón \emph{Decorator} (auditoría de pagos) y el
    patrón \emph{Adapter} (importación de Excel) se comporten según lo especificado.
    \item \textbf{Validar middlewares de seguridad:} comprobar el control de acceso por rol,
    la protección del último administrador activo y la expiración de sesión por
    inactividad.
    \item \textbf{Validar la integración con la capa de persistencia:} verificar que las
    operaciones sobre MongoDB (a través de \texttt{mongodb-memory-server}) sean correctas,
    aisladas entre pruebas y no dependan de una base de datos de producción.
\end{enumerate}

\clearpage

% ============================================================================
\section{Planteamiento}
% ============================================================================

\subsection{Alcance de las Pruebas}

Las pruebas unitarias e de integración ligera cubren:

\begin{itemize}[nosep]
    \item \textbf{Objetos de valor y modelo de dominio:} validación de datos, invariantes y
    métodos de negocio (\texttt{Cedula}, \texttt{Telefono}, \texttt{Usuario}).
    \item \textbf{Servicios de negocio:} reglas de negocio, validaciones, manejo de errores
    y orquestación transaccional (\texttt{AuthService}, \texttt{CopropietarioService},
    \texttt{PagoService}).
    \item \textbf{Patrones de diseño:} \texttt{EcuandorianMora12Strategy} /
    \texttt{CalculadorMoraContext} (Strategy), \texttt{AuditPaymentServiceDecorator}
    (Decorator) y \texttt{ExcelCopropietarioAdapter} (Adapter).
    \item \textbf{Middlewares de presentación:} \texttt{verificarSesion},
    \texttt{permitirSolo} y \texttt{verificarUltimoAdministrador}.
    \item \textbf{Controladores HTTP (integración):} \texttt{copropietarioController} mediante
    \texttt{supertest}, verificando el contrato completo petición-respuesta contra una base
    de datos en memoria.
\end{itemize}

Explícitamente \textbf{no} forman parte del alcance de este plan de pruebas automatizadas
(se mencionan como trabajo futuro en la Sección~7.4):

\begin{itemize}[nosep]
    \item La interfaz visual (HTML/CSS/JavaScript del lado cliente en
    \texttt{dashboard.html} y \texttt{login.html}).
    \item La generación de los binarios PDF/Excel (\texttt{generarPdfLista},
    \texttt{generarPdfReportePagos}, \texttt{generarPdfRecibo}, \texttt{generarExcelLista}),
    dado que su verificación exige comparación de contenido binario, no de valores lógicos.
    \item Pruebas de carga, rendimiento o end-to-end de navegador.
\end{itemize}

\subsection{Enfoque de Testing}

\textbf{Metodología:} desarrollo iterativo con pruebas de regresión (TDD adaptado) —
cada corrección de comportamiento de negocio (por ejemplo, el recálculo del recargo por
mora, ver Sección~6.2) fue acompañada de la actualización de su prueba correspondiente
antes de dar por cerrado el cambio.\\
\textbf{Tipo de pruebas:} unitarias puras (con \emph{mocks}) e integración ligera contra
una instancia de MongoDB en memoria.\\
\textbf{Cobertura objetivo:} no existe un umbral de cobertura de código configurado en
\texttt{jest.config.js}; en su lugar se verificó manualmente que cada requisito funcional
del backlog (\texttt{REQ001}–\texttt{REQ013}) cuente con al menos un caso de prueba
explícito, documentado en la Sección~4 y consolidado en la Tabla~\ref{tab:metricas-req}.

\subsection{Estrategia de Mocking}

AliGest combina \textbf{dos} estrategias de prueba complementarias, elegidas según el
tipo de componente:

\begin{enumerate}[nosep]
    \item \textbf{Mocks explícitos de Jest (\texttt{jest.mock()}):} utilizados quando el
    componente bajo prueba debe aislarse por completo de la capa de datos —middlewares de
    presentación y el decorador de auditoría—, ya que su responsabilidad es puramente de
    control de flujo, no de persistencia.
    \item \textbf{Base de datos en memoria (\texttt{mongodb-memory-server}):} utilizada para
    los servicios de negocio (\texttt{AuthService}, \texttt{CopropietarioService},
    \texttt{PagoService}) y para las pruebas de integración de controladores
    (\texttt{copropietarioController.test.js}), porque su lógica depende de consultas
    encadenadas reales (\texttt{populate}, unicidad de \texttt{cedula}/\texttt{casa},
    amortización FIFO sobre colecciones relacionadas) que un mock manual reproduciría de
    forma incompleta o frágil.
\end{enumerate}

\begin{lstlisting}[caption={Aislamiento con \texttt{jest.mock()} en \texttt{authMiddleware.test.js}}]
const { verificarSesion, permitirSolo, verificarUltimoAdministrador } =
    require('../../../src/presentation/middlewares/authMiddleware');
const usuarioRepository = require('../../../src/data/repositories/usuarioRepository');

jest.mock('../../../src/data/repositories/usuarioRepository');
\end{lstlisting}

\begin{lstlisting}[caption={Integración real con \texttt{mongodb-memory-server} en \texttt{tests/setup.js}}]
const { MongoMemoryServer } = require('mongodb-memory-server');
const mongoose = require('mongoose');
const db = require('../config/database');

let mongoServer;

beforeAll(async () => {
    mongoServer = await MongoMemoryServer.create();
    const uri = mongoServer.getUri();
    process.env.MONGODB_URI = uri;
    await db.conectarBD();
});

afterAll(async () => {
    await db.desconectarBD();
    if (mongoServer) {
        await mongoServer.stop();
    }
});

beforeEach(async () => {
    // Limpiar todas las colecciones antes de cada prueba (aislamiento entre casos)
    await db.limpiarTablas();
});
\end{lstlisting}

\subsection{Plan de Ejecución de Pruebas (Hoja de Ruta)}

Antes de implementar los casos de prueba descritos en la Sección~4, se definió la
siguiente hoja de ruta: un mapa de correspondencia entre cada requisito funcional del
backlog, el componente de código que lo materializa y el archivo de prueba planificado
para verificarlo. Esta planificación —reconstruida aquí a partir del trabajo ya
ejecutado— resume el orden y el alcance que el lector debe esperar encontrar en el
resto del informe.

\begin{table}[H]
\centering
\caption{Hoja de ruta de pruebas planificadas por requisito}
\label{tab:hoja-ruta}
\begin{tabular}{clp{3.7cm}p{3.6cm}l}
\toprule
\textbf{\#} & \textbf{REQ} & \textbf{Componente / Capa} & \textbf{Archivo de prueba planificado} & \textbf{Tipo} \\
\midrule
1  & REQ001 & \texttt{AuthService} + \texttt{authMiddleware} + \texttt{sessionRepository} & \texttt{AuthService.test.js} / \texttt{authMiddleware.test.js} / \texttt{authController.test.js} / \texttt{sessionRepository.test.js} & Integración + Unitaria \\
2  & REQ002 & \texttt{AuthService} (recuperación) & \texttt{AuthService.test.js} & Integración \\
3  & REQ003 & \texttt{authMiddleware} + \texttt{authController} (perfiles) & \texttt{authMiddleware.test.js} / \texttt{authController.test.js} / \texttt{sessionRepository.test.js} & Unitaria + Integración \\
4  & REQ004 & \texttt{CopropietarioService} + \texttt{excelAdapter} & \texttt{CopropietarioService.test.js} / \texttt{excelAdapter.test.js} / \texttt{ExcelImportSample.test.js} & Integración + Unitaria \\
5  & REQ005 & \texttt{copropietarioController} & \texttt{copropietarioController.test.js} & Integración HTTP \\
6  & REQ006 & \texttt{PagoService} (reportes) & \emph{Cobertura indirecta (Sección 7.4)} & Integración \\
7  & REQ007 & \texttt{CopropietarioService} (modificar) & \texttt{CopropietarioService.test.js} & Integración \\
8  & REQ008 & \texttt{CopropietarioService} + \texttt{copropietarioController} & \texttt{CopropietarioService.test.js} / \texttt{copropietarioController.test.js} & Integración \\
9  & REQ009 & \texttt{PagoService} + \texttt{pagoController} (reportar pago) & \texttt{PagoService.test.js} / \texttt{pagoController.test.js} & Integración \\
10 & REQ010 & \texttt{PagoService} + \texttt{auditDecorator} + \texttt{pagoController} & \texttt{PagoService.test.js} / \texttt{auditDecorator.test.js} / \texttt{pagoController.test.js} & Integración + Unitaria \\
11 & REQ011 & \texttt{moraStrategy} + \texttt{PagoService} (expensas) & \texttt{moraStrategy.test.js} / \texttt{PagoService.test.js} & Unitaria + Integración \\
12 & REQ012 & \texttt{Pago} + \texttt{pagoController} (recibo inmutable) & \texttt{pagoController.test.js} & Integración \\
13 & REQ013 & \texttt{NotificationService} & \texttt{NotificationService.test.js} & Integración \\
-- & Transversal & \texttt{Usuario}, \texttt{Cedula}, \texttt{Telefono} & \texttt{Usuario.test.js} / \texttt{Cedula.test.js} / \texttt{Telefono.test.js} & Unitaria \\
-- & Regresión & Vista \texttt{dashboard.html} & \texttt{dashboard.test.js} & Unitaria \\
\bottomrule
\end{tabular}
\end{table}

\begin{table}[H]
\centering
\caption{Orden de ejecución sugerido por fase de la arquitectura}
\label{tab:fases-plan}
\begin{tabular}{clp{8.5cm}}
\toprule
\textbf{Fase} & \textbf{Capa} & \textbf{Alcance} \\
\midrule
1 & Dominio y Objetos de Valor & Validar invariantes de \texttt{Cedula}, \texttt{Telefono} y \texttt{Usuario} antes de probar cualquier servicio que dependa de ellos. \\
2 & Patrones de Diseño & Verificar \texttt{Strategy} (mora), \texttt{Adapter} (Excel) y \texttt{Decorator} (auditoría) de forma aislada. \\
3 & Servicios de Negocio & Probar \texttt{AuthService}, \texttt{CopropietarioService} y \texttt{PagoService} contra \texttt{mongodb-memory-server}, una vez validadas sus dependencias de dominio y patrones. \\
4 & Presentación & Cerrar con los \emph{middlewares} de seguridad y las pruebas de integración HTTP sobre los controladores. \\
\bottomrule
\end{tabular}
\end{table}

Como se detalla en la Sección~4, de las 13 líneas planificadas, 12 se implementaron con
cobertura completa y 1 (REQ006, generación de reportes) quedó con cobertura parcial
—la restricción por rol se ejercita de forma indirecta mediante las pruebas de
integración de controladores, pero \texttt{obtenerReportePagos}/\texttt{generarPdfReportePagos}
aún no tienen un caso dedicado—; esa brecha se documenta de forma explícita en la
Sección~7.4 (Recomendaciones Futuras) en lugar de omitirse.
\clearpage

% ============================================================================
\section{Herramientas y Entorno}
% ============================================================================

\subsection{Framework de Testing}

\textbf{Jest 30.4.2} — framework de pruebas de JavaScript, ejecutado sobre Node.js en
modo CommonJS (\texttt{require}/\texttt{module.exports}), sin necesidad de transpilación.

\subsection{Entorno de Pruebas}

\begin{itemize}[nosep]
    \item \textbf{Node.js:} entorno de ejecución \texttt{testEnvironment: 'node'} (no se
    requiere DOM, ya que la capa de presentación validada es HTTP, no navegador).
    \item \textbf{Base de datos:} \texttt{mongodb-memory-server 11.2.0}, instancia efímera
    de MongoDB que se crea antes de la suite y se destruye al finalizar.
    \item \textbf{Peticiones HTTP simuladas:} \texttt{supertest 7.2.2}, para ejercer la
    aplicación Express completa (\texttt{app.js}) sin necesidad de abrir un puerto real.
    \item \textbf{Ejecución:} \texttt{jest --runInBand} (definido en
    \texttt{package.json}), en serie para evitar condiciones de carrera sobre la misma
    base de datos en memoria compartida entre archivos de prueba.
\end{itemize}

\subsection{Estructura de Archivos}

\begin{lstlisting}[language={},caption={Estructura del directorio \texttt{tests/}}]
tests/
├── setup.js
├── core/
│   ├── business/
│   │   ├── patterns/
│   │   │   ├── auditDecorator.test.js      (2 tests)
│   │   │   ├── excelAdapter.test.js        (4 tests)
│   │   │   └── moraStrategy.test.js        (4 tests)
│   │   └── services/
│   │       ├── AuthService.test.js         (9 tests)
│   │       ├── CopropietarioService.test.js(7 tests)
│   │       ├── ExcelImportSample.test.js   (1 test)
│   │       ├── NotificationService.test.js (3 tests)
│   │       └── PagoService.test.js         (10 tests)
│   └── domain/
│       ├── models/
│       │   └── Usuario.test.js             (7 tests)
│       └── valueObjects/
│           ├── Cedula.test.js              (7 tests)
│           └── Telefono.test.js            (7 tests)
├── data/
│   └── repositories/
│       └── sessionRepository.test.js       (3 tests)
└── presentation/
    ├── controllers/
    │   ├── authController.test.js          (4 tests)
    │   ├── copropietarioController.test.js (5 tests)
    │   └── pagoController.test.js          (4 tests)
    ├── middlewares/
    │   └── authMiddleware.test.js          (6 tests)
    └── views/
        └── dashboard.test.js               (3 tests)
\end{lstlisting}

\clearpage

% ============================================================================
\section{Implementación por Requisito}
% ============================================================================

Esta sección documenta, para cada requisito funcional (\texttt{REQ001}–\texttt{REQ013},
según la numeración consolidada en el backlog del proyecto), el código relevante que
implementa la funcionalidad, el archivo de prueba correspondiente y los resultados
obtenidos.

% ----------------------------------------------------------------------------
\subsection{REQ001: Gestión de Acceso y Registro de Usuarios}
% ----------------------------------------------------------------------------
\textbf{Archivos:} \texttt{AuthService.test.js} (login, registro, bloqueo y desbloqueo),
\texttt{authMiddleware.test.js} (\texttt{verificarSesion}/\texttt{permitirSolo}),
\texttt{authController.test.js} (registro público bloqueado) y
\texttt{sessionRepository.test.js} (sesión aleatoria y expiración).\\
\textbf{Total de pruebas:} 12 — \textbf{Estado:} 12/12 PASANDO (100\,\%).

\subsubsection*{Cómo se implementó}

El registro y el inicio de sesión se validan primero contra las reglas de dominio del
objeto \texttt{Usuario} y luego contra el repositorio de persistencia:

\begin{lstlisting}[caption={\texttt{AuthService.login} — validación de credenciales y bloqueo temporal (\texttt{src/core/business/services/AuthService.js})}]
async login(username, password) {
    if (!username || !password) {
        throw new Error("El usuario y la contraseña son obligatorios.");
    }

    const user = await usuarioRepository.findByUsername(username);
    if (!user) {
        throw new Error("Nombre de usuario o contraseña incorrectos. Intente nuevamente.");
    }

    const domainUser = new Usuario(user);
    if (domainUser.esBloqueado()) {
        const limite = new Date(user.lockout_until);
        const minRestantes = Math.ceil((limite - new Date()) / (1000 * 60));
        throw new Error(`Cuenta bloqueada temporalmente por múltiples intentos fallidos. Intente en ${minRestantes} minutos o recupere su contraseña.`);
    }

    const valid = await bcrypt.compare(password, user.password);
    if (!valid) {
        await usuarioRepository.incrementFailedAttempts(username);
        const updatedUser = await usuarioRepository.findByUsername(username);
        if (updatedUser.failed_attempts >= 3) {
            const lockoutTime = new Date(Date.now() + 15 * 60 * 1000).toISOString();
            await usuarioRepository.lockoutUser(username, lockoutTime);
            throw new Error("Cuenta bloqueada temporalmente por múltiples intentos fallidos. Intente en 15 minutos o recupere su contraseña.");
        }
        const intentosRestantes = 3 - updatedUser.failed_attempts;
        throw new Error(`Nombre de usuario o contraseña incorrectos. Le quedan ${intentosRestantes} intentos antes de bloquear la cuenta.`);
    }

    await usuarioRepository.resetFailedAttempts(username);
    return { id: user.id, username: user.username, email: user.email, role: user.role, mustChangePassword: user.must_change_password === 1 };
}
\end{lstlisting}

El middleware \texttt{verificarSesion} controla, además del rol, la expiración de sesión
por inactividad (RNF-01, ver Sección~4.11 y el detalle del mecanismo en el Anexo~C):

\begin{lstlisting}[caption={Prueba de bloqueo por 3 intentos fallidos (\texttt{AuthService.test.js})}]
test('Debería bloquear la cuenta por 15 minutos en el tercer intento fallido de inicio de sesión', async () => {
    const username = 'bloqueable';
    const passwordCorrecta = 'Password123!';
    await authService.registrarUsuarioPublico(username, 'bloqueo@correo.com', passwordCorrecta);

    await expect(authService.login(username, 'Incorrecta')).rejects.toThrow("Le quedan 2 intentos");
    await expect(authService.login(username, 'Incorrecta')).rejects.toThrow("Le quedan 1 intentos");
    await expect(authService.login(username, 'Incorrecta'))
        .rejects.toThrow("Cuenta bloqueada temporalmente por múltiples intentos fallidos");

    await expect(authService.login(username, passwordCorrecta))
        .rejects.toThrow("Cuenta bloqueada temporalmente");
});
\end{lstlisting}

\subsubsection*{Resultados Obtenidos}
\begin{itemize}[nosep]
    \item Registro público con contraseña cifrada (\texttt{bcrypt}, factor 12) — verificado.
    \item Rechazo de usuario/correo duplicado — verificado.
    \item Inicio de sesión válido con carga de permisos por rol — verificado.
    \item Bloqueo progresivo (2 intentos restantes $\rightarrow$ 1 $\rightarrow$ bloqueo de
    15 minutos) y persistencia del bloqueo ante credenciales correctas — verificado.
    \item Autorización por encabezados \texttt{x-user-role}/\texttt{x-user-id} y rechazo de
    roles inválidos (401/403) — verificado.
    \item Restricción de rutas por rol mediante \texttt{permitirSolo(...)} — verificado.
\end{itemize}

% ----------------------------------------------------------------------------
\subsection{REQ002: Recuperación de Contraseña}
% ----------------------------------------------------------------------------
\textbf{Archivo:} \texttt{tests/core/business/services/AuthService.test.js} (4 casos:
generación del OTP, restablecimiento, expiración del código e invalidación tras 3 intentos).\\
\textbf{Total de pruebas:} 4 — \textbf{Estado:} 4/4 PASANDO (100\,\%).

\subsubsection*{Cómo se implementó}

\begin{lstlisting}[caption={Generación y expiración del código OTP (\texttt{AuthService.js})}]
async generarCodigoRecuperacion(email) {
    if (!email) {
        throw new Error("El correo electrónico es obligatorio para recuperar contraseña.");
    }
    const user = await usuarioRepository.findByEmail(email);
    if (!user) {
        throw new Error("El correo electrónico ingresado no está registrado en el sistema. Contacte al administrador.");
    }
    const codigo = Math.floor(100000 + Math.random() * 900000).toString();
    const fechaExpiracion = new Date(Date.now() + 10 * 60 * 1000); // vigencia de 10 minutos
    await usuarioRepository.updateRecoveryCode(user.username, codigo, fechaExpiracion.toISOString());
    // ... simulación de despacho por SMTP (consola) ...
    return { mensaje: "Código enviado correctamente.", codigoSimulado: codigo, username: user.username };
}
\end{lstlisting}

\begin{lstlisting}[caption={Prueba de expiración del código OTP (\texttt{AuthService.test.js})}]
test('Debería fallar al cambiar clave si el código OTP ha expirado', async () => {
    const username = 'expirable';
    const email = 'expirable@correo.com';
    await authService.registrarUsuarioPublico(username, email, 'OldPassword123!');

    await authService.generarCodigoRecuperacion(email);
    const expiracionPasada = new Date(Date.now() - 10 * 1000).toISOString();
    const user = await usuarioRepository.findByUsername(username);
    await usuarioRepository.updateRecoveryCode(username, user.recovery_code, expiracionPasada);

    await expect(authService.recuperarContrasena(username, user.recovery_code, 'NewPassword123!'))
        .rejects.toThrow("El código de verificación ha expirado. Solicite uno nuevo.");
});
\end{lstlisting}

\subsubsection*{Resultados Obtenidos}
\begin{itemize}[nosep]
    \item Generación de código de 6 dígitos con vigencia de 10 minutos — verificado.
    \item Restablecimiento exitoso invalida la contraseña anterior — verificado.
    \item Rechazo de código expirado con mensaje explícito — verificado.
\end{itemize}

% ----------------------------------------------------------------------------
\subsection{REQ003: Configuración de Perfiles}
% ----------------------------------------------------------------------------
\textbf{Archivos:} \texttt{authMiddleware.test.js} (protección del último administrador),
\texttt{authController.test.js} (promoción de perfil, listado con nombre/casa y bloqueo de
degradación del único administrador) y \texttt{sessionRepository.test.js} (invalidación de
sesiones al cambiar el perfil).\\
\textbf{Total de pruebas:} 6 — \textbf{Estado:} 6/6 PASANDO (100\,\%).

\subsubsection*{Cómo se implementó}

\begin{lstlisting}[caption={\texttt{verificarUltimoAdministrador} (\texttt{authMiddleware.js})}]
async function verificarUltimoAdministrador(req, res, next) {
    const { id } = req.params;
    if (!id) return next();

    const usuarioTarget = await usuarioRepository.findById(id);
    if (!usuarioTarget) {
        return res.status(404).json({ error: "El usuario especificado no existe en el sistema." });
    }

    if (usuarioTarget.role === 'ADMINISTRADOR') {
        const numAdmins = await usuarioRepository.countAdministrators();
        if (numAdmins <= 1) {
            return res.status(400).json({
                error: "Operación de seguridad denegada: No se puede eliminar o desactivar al último administrador del sistema."
            });
        }
    }
    next();
}
\end{lstlisting}

\begin{lstlisting}[caption={Prueba del bloqueo al último administrador activo (\texttt{authMiddleware.test.js})}]
test('Debería bloquear la acción (400) si se intenta eliminar o desactivar al último administrador activo del sistema', async () => {
    mockReq.params.id = 'ultimo-admin';
    usuarioRepository.findById.mockResolvedValue({ id: 'ultimo-admin', role: 'ADMINISTRADOR' });
    usuarioRepository.countAdministrators.mockResolvedValue(1); // Solo queda 1 admin

    await verificarUltimoAdministrador(mockReq, mockRes, nextFunction);

    expect(mockRes.status).toHaveBeenCalledWith(400);
    expect(mockRes.json).toHaveBeenCalledWith(expect.objectContaining({
        error: expect.stringContaining("No se puede eliminar o desactivar al último administrador")
    }));
    expect(nextFunction).not.toHaveBeenCalled();
});
\end{lstlisting}

\subsubsection*{Resultados Obtenidos}
\begin{itemize}[nosep]
    \item Modificación de perfiles de usuarios distintos de ADMINISTRADOR sin restricción — verificado.
    \item Degradación de un ADMINISTRADOR permitida cuando existen otros administradores activos — verificado.
    \item Bloqueo (400) cuando la operación dejaría al sistema sin ningún administrador — verificado.
\end{itemize}

% ----------------------------------------------------------------------------
\subsection{REQ004: Importación de Copropietarios}
% ----------------------------------------------------------------------------
\textbf{Archivos:} \texttt{CopropietarioService.test.js} (3 casos: creación manual e
importación masiva), \texttt{excelAdapter.test.js} (4 casos) y
\texttt{ExcelImportSample.test.js} (importación aislada de 60 filas con creación de cuentas).\\
\textbf{Total de pruebas:} 8 — \textbf{Estado:} 8/8 PASANDO (100\,\%).

\subsubsection*{Cómo se implementó}

El patrón \emph{Adapter} tolera variaciones de encabezado del archivo Excel de origen:

\begin{lstlisting}[caption={\texttt{ExcelCopropietarioAdapter.adaptRow} (\texttt{src/core/business/patterns/excelAdapter.js})}]
static adaptRow(externalRow) {
    if (!externalRow) {
        throw new Error("Fila de hoja de cálculo nula o indefinida.");
    }
    const cedula = externalRow["Cédula"] || externalRow["Cedula"] || externalRow["ID_Nacional"] || "";
    const nombre = externalRow["Nombre Completo"] || externalRow["Nombre"] || "";
    const casa = externalRow["Número Casa"] || externalRow["Casa"] || "";
    const telefono = externalRow["Teléfono"] || externalRow["Telefono"] || externalRow["Celular"] || "";
    const email = externalRow["Correo Electrónico"] || externalRow["Correo"] || externalRow["Email"] || "";
    const saldo = parseFloat(externalRow["Saldo Inicial"] || externalRow["Saldo"] || 0.0);

    return {
        cedula: String(cedula).trim(), nombre: String(nombre).trim(), casa: String(casa).trim(),
        telefono: String(telefono).trim(), email: String(email).trim(),
        saldo: isNaN(saldo) ? 0.0 : saldo
    };
}
\end{lstlisting}

La importación masiva es \textbf{atómica}: valida todo el lote antes de persistir cualquier
registro, garantizando que un solo error cancele la operación completa (REQ004):

\begin{lstlisting}[caption={Prueba de atomicidad de la importación masiva (\texttt{CopropietarioService.test.js})}]
test('Debería abortar toda la importación masiva si existe una sola fila inválida (Atomicidad)', async () => {
    const filas = [
        { "Cédula": "1721151247", "Nombre Completo": "Juan Perez", "Número Casa": "Villa 1", "Correo": "juan@correo.com" },
        { "Cédula": "1721151241", "Nombre Completo": "Fila Invalida", "Número Casa": "Villa 2", "Correo": "invalido" }
    ];
    const buffer = crearExcelMock(filas);

    await expect(copropietarioService.importarMasivoDesdeExcel(buffer))
        .rejects.toThrow("Importación cancelada por errores de validación");

    const todos = await copropietarioService.obtenerTodos();
    expect(todos).toHaveLength(0); // Nada se persistió
});
\end{lstlisting}

\subsubsection*{Resultados Obtenidos}
\begin{itemize}[nosep]
    \item Registro manual individual con generación automática de usuario/contraseña
    temporal (\texttt{must\_change\_password = 1}) — verificado.
    \item Importación masiva exitosa de 2 filas con generación de \texttt{username} a
    partir del número de casa — verificado.
    \item Cancelación total del lote ante una fila con cédula/correo inválidos — verificado.
    \item Mapeo tolerante de encabezados en español, alternativos y valores de saldo no
    numéricos — verificado (\texttt{excelAdapter.test.js}).
\end{itemize}

% ----------------------------------------------------------------------------
\subsection{REQ005: Consulta de Copropietarios}
% ----------------------------------------------------------------------------
\textbf{Archivo:} \texttt{tests/presentation/controllers/copropietarioController.test.js}
(3 casos, integración HTTP con \texttt{supertest}: listado autenticado, filtro por nombre y
restricción de la nómina/\texttt{/me} para el rol Copropietario).\\
\textbf{Total de pruebas:} 3 — \textbf{Estado:} 3/3 PASANDO (100\,\%).

\subsubsection*{Cómo se implementó}

\begin{lstlisting}[caption={Filtro por nombre en \texttt{GET /api/copropietarios} (\texttt{src/presentation/controllers/copropietarioController.js})}]
router.get('/', verificarSesion, async (req, res) => {
    const { nombre, casa, casaInicio, casaFin, estadoCuenta } = req.query;
    let lista = await copropietarioService.obtenerTodos();

    if (nombre) {
        const queryName = nombre.toLowerCase().trim();
        lista = lista.filter(c => c.nombre.toLowerCase().includes(queryName));
    }
    // ... filtros adicionales por casa, rango de villas y estado de cuenta ...
    res.status(200).json(lista);
});
\end{lstlisting}

\begin{lstlisting}[caption={Prueba de integración con \texttt{supertest} (\texttt{copropietarioController.test.js})}]
test('GET /api/copropietarios - Debería filtrar residentes por nombre', async () => {
    await Copropietario.create({
        cedula: '0921151247', nombre: 'Maria Gomez', casa: 'Villa 20',
        telefono: '+593987654322', email: 'maria@correo.com', saldo: 0.0
    });

    const response = await request(app)
        .get('/api/copropietarios?nombre=Maria')
        .set(adminTokenHeaders);

    expect(response.status).toBe(200);
    expect(response.body).toHaveLength(1);
    expect(response.body[0].nombre).toBe('Maria Gomez');
});
\end{lstlisting}

\subsubsection*{Resultados Obtenidos}
\begin{itemize}[nosep]
    \item Listado completo de copropietarios autenticados (rol Administrador) — verificado.
    \item Filtro por nombre insensible a mayúsculas — verificado.
\end{itemize}

\subsubsection*{Nota de Cobertura}
Los filtros por rango de villas, estado de cuenta (\texttt{saldado}/\texttt{deudor}) y el
ordenamiento de columnas se implementan en el mismo controlador, pero no cuentan con un
caso de prueba automatizado dedicado; fueron verificados manualmente durante el desarrollo.
Se recomienda ampliar \texttt{copropietarioController.test.js} con estos escenarios (ver
Sección~7.4).

% ----------------------------------------------------------------------------
\subsection{REQ006: Generación de Reportes}
% ----------------------------------------------------------------------------
\textbf{Total de pruebas:} 0 — \textbf{Estado:} SIN COBERTURA AUTOMATIZADA DEDICADA.

\subsubsection*{Cómo se implementó}

\texttt{PagoService.obtenerReportePagos} restringe automáticamente el conjunto de
resultados según el rol del usuario autenticado:

\begin{lstlisting}[caption={Restricción de reportes por rol (\texttt{src/core/business/services/PagoService.js})}]
async obtenerReportePagos(filtros, userContext) {
    const role = userContext?.role;
    let copropietarioId = null;

    if (role === 'COPROPIETARIO') {
        const copro = await copropietarioRepository.findByUsuarioId(userContext.id);
        if (!copro) return [];
        copropietarioId = copro.id;
    }

    return await pagoRepository.getReportePagos({
        fechaInicio: filtros.fechaInicio || null,
        fechaFin: filtros.fechaFin || null,
        periodo: filtros.periodo || null,
        estado: filtros.estado || null,
        copropietarioId
    });
}
\end{lstlisting}

\subsubsection*{Resultados Obtenidos y Nota de Cobertura}
La lógica de restricción por rol reutiliza el mismo mecanismo ya validado indirectamente
en \texttt{PagoService.test.js} (búsqueda de copropietario por \texttt{usuarioId}), pero no
existe, al cierre de este informe, un caso de prueba que invoque
\texttt{obtenerReportePagos} ni \texttt{generarPdfReportePagos} directamente. Se identifica
como una brecha de cobertura y se registra como recomendación prioritaria en la
Sección~7.4.

% ----------------------------------------------------------------------------
\subsection{REQ007: Modificación de Copropietarios}
% ----------------------------------------------------------------------------
\textbf{Archivo:} \texttt{tests/core/business/services/CopropietarioService.test.js}
(3 casos).\\
\textbf{Total de pruebas:} 3 — \textbf{Estado:} 3/3 PASANDO (100\,\%).

\subsubsection*{Cómo se implementó}

Además de la actualización estándar de datos de contacto, \texttt{actualizarDatos} permite
el cambio de cédula (representante nuevo de la vivienda) regenerando automáticamente la
contraseña temporal de acceso:

\begin{lstlisting}[caption={Regeneración de contraseña ante cambio de cédula (\texttt{CopropietarioService.js})}]
if (datos.cedula) {
    const cedulaSolicitada = new Cedula(datos.cedula).valor;
    if (cedulaSolicitada !== copropietario.cedula) {
        const existeCedula = await copropietarioRepository.findByCedula(cedulaSolicitada);
        if (existeCedula) {
            throw new Error(`La cédula ${cedulaSolicitada} ya está registrada en el sistema.`);
        }
        nuevaCedula = cedulaSolicitada;

        if (copropietario.usuario_id) {
            passwordTemporal = `Temp-${Math.floor(1000 + Math.random() * 9000)}!`;
            const hashed = await bcrypt.hash(passwordTemporal, 12);
            await usuarioRepository.updatePasswordAndForceMustChange(copropietario.usuario_id, hashed);
        }
    }
}
\end{lstlisting}

\begin{lstlisting}[caption={Prueba del cambio de representante (\texttt{CopropietarioService.test.js})}]
test('RF-2.4: Debería permitir cambiar la cédula (nuevo representante) y regenerar la contraseña temporal', async () => {
    await copropietarioService.crearCopropietario({
        cedula: '1721151247', nombre: 'Juan Original', casa: 'Villa 2',
        telefono: '0987654321', email: 'juan2@correo.com'
    });

    const actualizado = await copropietarioService.actualizarDatos('1721151247', {
        cedula: '3000000012', nombre: 'Pedro Nuevo Representante'
    });

    expect(actualizado.cedula).toBe('3000000012');
    expect(actualizado.passwordTemporal).toEqual(expect.stringMatching(/^Temp-\d{4}!$/));

    const anterior = await copropietarioRepository.findByCedula('1721151247');
    expect(anterior).toBeNull();

    const nuevo = await copropietarioRepository.findByCedula('3000000012');
    const usuario = await usuarioRepository.findById(nuevo.usuario_id);
    expect(usuario.must_change_password).toBe(1);
});
\end{lstlisting}

\subsubsection*{Resultados Obtenidos}
\begin{itemize}[nosep]
    \item Actualización de nombre, villa, teléfono y correo — verificado.
    \item Cambio de cédula con regeneración de contraseña temporal y forzado de cambio en
    el primer ingreso — verificado.
    \item Rechazo del cambio si la nueva cédula ya pertenece a otro copropietario — verificado.
\end{itemize}

% ----------------------------------------------------------------------------
\subsection{REQ008: Eliminación de Copropietarios}
% ----------------------------------------------------------------------------
\textbf{Archivos:} \texttt{CopropietarioService.test.js} (1 caso) y
\texttt{copropietarioController.test.js} (2 casos, integración HTTP).\\
\textbf{Total de pruebas:} 3 — \textbf{Estado:} 3/3 PASANDO (100\,\%).

\subsubsection*{Cómo se implementó}

El controlador exige confirmación explícita cuando el copropietario tiene historial de
pagos, y bloquea por completo la eliminación si tiene deudas pendientes:

\begin{lstlisting}[caption={Confirmación por historial de pagos (\texttt{copropietarioController.js})}]
router.delete('/:id', verificarSesion, permitirSolo('ADMINISTRADOR'), async (req, res) => {
    const id = req.params.id;
    const { confirmar } = req.query;

    const deudasPendientes = await pagoRepository.countPendingDeudas(id);
    if (deudasPendientes > 0) {
        return res.status(400).json({ error: "No se permite eliminar un residente con deudas pendientes." });
    }

    const tienePagos = await pagoRepository.countPayments(id);
    if (tienePagos > 0 && confirmar !== 'true') {
        return res.status(202).json({
            advertencia: "HISTORIAL_ACTIVO",
            mensaje: "ADVERTENCIA: Este copropietario tiene historial de pagos registrado. Eliminar este registro puede afectar reportes históricos y estadísticas. ¿Desea continuar de todas formas?"
        });
    }
    const result = await copropietarioService.eliminarCopropietario(id);
    res.status(200).json(result);
});
\end{lstlisting}

\begin{lstlisting}[caption={Prueba de la advertencia 202 por historial activo (\texttt{copropietarioController.test.js})}]
test('DELETE /api/copropietarios/:id - Debería denegar con 202 y advertencia si tiene historial de pagos sin confirmar', async () => {
    await Pago.create({
        comprobante_id: 'DEP-777', copropietario_id: coproId, monto_pagado: 100.0,
        estado: 'APROBADO', fecha_registro: new Date().toISOString()
    });

    const response = await request(app)
        .delete(`/api/copropietarios/${coproId}`)
        .set(adminTokenHeaders);

    expect(response.status).toBe(202);
    expect(response.body.advertencia).toBe('HISTORIAL_ACTIVO');
});
\end{lstlisting}

\subsubsection*{Resultados Obtenidos}
\begin{itemize}[nosep]
    \item Bloqueo total de la eliminación si existen deudas pendientes — verificado.
    \item Advertencia (202) si existe historial de pagos sin confirmación explícita — verificado.
    \item Eliminación (borrado lógico del copropietario y físico del usuario) al confirmar
    con \texttt{?confirmar=true} — verificado.
\end{itemize}

% ----------------------------------------------------------------------------
\subsection{REQ009: Registro de Pago de Expensa}
% ----------------------------------------------------------------------------
\textbf{Archivos:} \texttt{PagoService.test.js} (2 casos: reporte pendiente y rechazo de
monto no positivo) y \texttt{pagoController.test.js} (2 casos de integración HTTP:
conservación de la imagen del comprobante y rechazo de un archivo que finge ser PNG).\\
\textbf{Total de pruebas:} 4 — \textbf{Estado:} 4/4 PASANDO (100\,\%).

\subsubsection*{Cómo se implementó}

\begin{lstlisting}[caption={\texttt{reportarPagoCopropietario} (\texttt{PagoService.js})}]
async reportarPagoCopropietario(usuarioId, comprobanteId, monto, metodo, periodo, comprobanteImg) {
    const copro = await copropietarioRepository.findByUsuarioId(usuarioId);
    if (!copro) {
        throw new Error("No se encontró el registro de copropietario para el usuario actual.");
    }

    const pagoEntidad = new Pago({ copropietarioId: copro.id, monto, comprobanteId, metodo, periodo, comprobanteImg });
    if (!pagoEntidad.esMontoValido()) {
        throw new Error("El monto reportado debe ser un número positivo.");
    }

    await pagoRepository.registrarPagoInmutable(
        pagoEntidad.comprobanteId, pagoEntidad.copropietarioId, pagoEntidad.monto,
        'PENDIENTE_VALIDACION', null, pagoEntidad.metodo, pagoEntidad.periodo, pagoEntidad.comprobanteImg
    );

    return { mensaje: "El pago fue reportado con éxito y está pendiente de verificación." };
}
\end{lstlisting}

\subsubsection*{Resultados Obtenidos}
\begin{itemize}[nosep]
    \item Reporte de pago queda en estado \texttt{PENDIENTE\_VALIDACION} con la referencia
    bancaria y el monto correctos — verificado.
    \item Rechazo de montos no positivos (\texttt{esMontoValido()}) — verificado.
\end{itemize}

\subsubsection*{Nota de Cobertura}
La validación de tipo de archivo (JPG/JPEG/PNG) y del límite de 5\,MB del comprobante
se implementa mediante \texttt{multer} directamente en la capa de controlador
(\texttt{pagoController.js}); al depender del \emph{middleware} de carga multipart, no se
ejerce en las pruebas de servicio aquí descritas y queda pendiente como prueba de
integración HTTP futura (Sección~7.4).

% ----------------------------------------------------------------------------
\subsection{REQ010: Validación de Pagos}
% ----------------------------------------------------------------------------
\textbf{Archivos:} \texttt{PagoService.test.js} (reserva atómica anti–doble aprobación y
rechazo con motivo), \texttt{auditDecorator.test.js} (2 casos del patrón Decorator) y
\texttt{pagoController.test.js} (rechazo vía API con el ObjectId del pago).\\
\textbf{Total de pruebas:} 5 — \textbf{Estado:} 5/5 PASANDO (100\,\%).

\subsubsection*{Cómo se implementó}

La aprobación de un pago se envuelve con el patrón \emph{Decorator} para registrar
auditoría sin acoplar esa responsabilidad a la lógica de negocio:

\begin{lstlisting}[caption={\texttt{AuditPaymentServiceDecorator} (\texttt{src/core/business/patterns/auditDecorator.js})}]
function AuditPaymentServiceDecorator(baseService) {
    return {
        async procesarPago(pagoId) {
            await pagoRepository.registrarAuditoria("INTENTO_APROBACION_PAGO", { pagoId, timestamp: new Date().toISOString() });
            try {
                const resultado = await baseService.validarAprobarPago(pagoId);
                await pagoRepository.registrarAuditoria("PAGO_APROBADO_EXITO", {
                    pagoId, reciboId: resultado.comprobante_id, monto: resultado.monto_aplicado,
                    sobrepago: resultado.sobrepago, timestamp: new Date().toISOString()
                });
                return resultado;
            } catch (err) {
                await pagoRepository.registrarAuditoria("APROBACION_PAGO_FALLIDA", { pagoId, error: err.message, timestamp: new Date().toISOString() });
                throw err;
            }
        }
    };
}
\end{lstlisting}

\begin{lstlisting}[caption={Prueba de auditoría en caso de fallo (\texttt{auditDecorator.test.js}, con \texttt{jest.mock})}]
test('Debería registrar auditoría de fallo si el servicio base lanza un error', async () => {
    const mockError = new Error("Saldo insuficiente");
    mockBaseService.validarAprobarPago.mockRejectedValue(mockError);
    pagoRepository.registrarAuditoria.mockResolvedValue('AUDIT-ID');

    await expect(decoratedService.procesarPago('PAGO-123')).rejects.toThrow("Saldo insuficiente");

    expect(pagoRepository.registrarAuditoria).toHaveBeenCalledTimes(2);
    expect(pagoRepository.registrarAuditoria).toHaveBeenNthCalledWith(
        2, "APROBACION_PAGO_FALLIDA",
        expect.objectContaining({ pagoId: 'PAGO-123', error: "Saldo insuficiente" })
    );
});
\end{lstlisting}

\subsubsection*{Resultados Obtenidos}
\begin{itemize}[nosep]
    \item Aprobación de pago con amortización FIFO, generación de recibo digital y
    auditoría de éxito (\texttt{INTENTO\_APROBACION\_PAGO} $\rightarrow$
    \texttt{PAGO\_APROBADO\_EXITO}) — verificado.
    \item Auditoría de fallo (\texttt{APROBACION\_PAGO\_FALLIDA}) cuando el servicio base
    lanza una excepción — verificado.
    \item Rechazo de pago con motivo obligatorio y notificación al residente — verificado.
\end{itemize}

% ----------------------------------------------------------------------------
\subsection{REQ011: Cálculo de Sanción por Mora}
% ----------------------------------------------------------------------------
\textbf{Archivos:} \texttt{tests/core/business/patterns/moraStrategy.test.js} (4 casos) y
\texttt{PagoService.test.js} (6 casos: amortización FIFO con mora, no reaplicación del 12\,\%,
generación de expensas mensuales, no duplicación por período, validación de formato y
vencimiento con cambio de año — infraestructura necesaria para que exista una deuda real
que amortizar).\\
\textbf{Total de pruebas:} 10 — \textbf{Estado:} 10/10 PASANDO (100\,\%).

\subsubsection*{Cómo se implementó}

El recargo por mora se implementa con el patrón \emph{Strategy}, como un
\textbf{recargo fijo del 12\,\%} aplicado una única vez por período vencido (no se
prorratea por días de atraso ni se acumula):

\begin{lstlisting}[caption={\texttt{EcuandorianMora12Strategy} (\texttt{src/core/business/patterns/moraStrategy.js})}]
class CalculadorMoraContext {
    constructor(strategy) { this.strategy = strategy; }
    calcular(monto, diasRetraso) { return this.strategy.calcularMora(monto, diasRetraso); }
}

class EcuandorianMora12Strategy {
    calcularMora(monto, diasRetraso) {
        if (diasRetraso <= 0) return 0;
        // Recargo fijo del 12% sobre la deuda vencida, aplicado una sola
        // vez por período (no se prorratea ni se acumula por días de atraso).
        const interes = monto * 0.12;
        return parseFloat(interes.toFixed(2));
    }
}
\end{lstlisting}

\begin{lstlisting}[caption={Prueba de invariancia del recargo respecto a los días de atraso (\texttt{moraStrategy.test.js})}]
test('El recargo del 12% no debe acumularse ni variar según los días de retraso transcurridos', () => {
    // El recargo se aplica una sola vez por período vencido (RF-3.3),
    // por lo que el resultado es el mismo con 1 día o con 200 días de mora.
    expect(context.calcular(100, 1)).toBe(12.00);
    expect(context.calcular(100, 200)).toBe(12.00);
});
\end{lstlisting}

Dado que la generación de la deuda mensual (\texttt{Deuda}) es un prerrequisito para que
el cálculo de mora tenga un caso real que amortizar, se añadió y probó
\texttt{generarExpensasMensuales}, que registra la expensa del período para todos los
copropietarios sin duplicar períodos ya generados, con vencimiento el día 5 del mes
siguiente:

\begin{lstlisting}[caption={\texttt{generarExpensasMensuales} (\texttt{PagoService.js})}]
function calcularVencimientoExpensa(periodo) {
    const [anioStr, mesStr] = periodo.split('-');
    let anio = parseInt(anioStr, 10);
    let mes = parseInt(mesStr, 10) + 1; // mes (1-12) del período + 1 = mes siguiente
    if (mes > 12) { mes = 1; anio += 1; }
    return `${anio}-${String(mes).padStart(2, '0')}-05`;
}

async generarExpensasMensuales(periodo, montoBase) {
    if (!periodo || !/^\d{4}-(0[1-9]|1[0-2])$/.test(periodo)) {
        throw new Error("El período debe tener el formato YYYY-MM.");
    }
    const monto = parseFloat(montoBase);
    const fechaVencimiento = calcularVencimientoExpensa(periodo);
    const copropietarios = await copropietarioRepository.findAll();

    let generadas = 0, omitidas = 0;
    for (const copro of copropietarios) {
        const existente = await pagoRepository.findDeudaPorPeriodo(copro.id, periodo);
        if (existente) { omitidas++; continue; }
        await pagoRepository.createDeuda(copro.id, periodo, monto, 'PENDIENTE', fechaVencimiento);
        generadas++;
    }
    return { mensaje: `Expensas del período ${periodo} generadas correctamente.`, periodo, montoBase: monto, fechaVencimiento, generadas, omitidas, totalCopropietarios: copropietarios.length };
}
\end{lstlisting}

\begin{lstlisting}[caption={Amortización FIFO con recargo del 12\% sobre la deuda vencida (\texttt{PagoService.test.js})}]
test('Debería aplicar amortización FIFO e interés de mora del 12% al aprobar un pago', async () => {
    // Deuda 1: Vencida hace 10 días, monto $100. Deuda 2: Vigente, monto $100.
    const d1Id = await pagoRepository.createDeuda(copro.id, '2026-06', 100.00, 'PENDIENTE', '2026-07-10');
    const d2Id = await pagoRepository.createDeuda(copro.id, '2026-07', 100.00, 'PENDIENTE', '2026-07-30');

    await pagoService.reportarPagoCopropietario(usuarioId, 'TRANS-FIFO', 150.00, 'TRANSFERENCIA', '2026-07', 'img');
    const pagoId = (await pagoService.obtenerPagosPendientes())[0].id;

    const resultado = await pagoService.validarAprobarPago(pagoId, new Date('2026-07-20'));
    expect(resultado.estado_transaccion).toBe('APROBADO_E_INMUTABLE');

    // Deuda 1: 100 + (100*0.12) = 112.00 -> se cancela por completo (150 - 112 = 38 restante)
    // Deuda 2: sin mora (vigente) -> 100 - 38 = 62.00 pendiente
    const deudas = await pagoRepository.getDeudasByCopropietario(copro.id);
    expect(deudas.find(d => d.id === d1Id).estado).toBe('PAGADO');
    expect(deudas.find(d => d.id === d2Id).monto).toBe(62.00);
});
\end{lstlisting}

\subsubsection*{Resultados Obtenidos}
\begin{itemize}[nosep]
    \item Recargo cero para deudas no vencidas — verificado.
    \item Recargo fijo del 12\,\% sobre el monto adeudado, sin prorrateo diario — verificado.
    \item Invariancia del recargo frente a distintos períodos de atraso (1 día vs.\ 200
    días) — verificado.
    \item Amortización FIFO completa: cancelación total de la deuda más antigua y aplicación
    del remanente a la siguiente — verificado.
    \item Generación de expensas mensuales sin duplicar el período y con vencimiento
    correcto (incluido el cambio de año en diciembre $\rightarrow$ enero) — verificado.
\end{itemize}

% ----------------------------------------------------------------------------
\subsection{REQ012: Generación de Comprobante de Pago}
% ----------------------------------------------------------------------------
\textbf{Archivo:} \texttt{tests/presentation/controllers/pagoController.test.js} (1 caso de
integración HTTP).\\
\textbf{Total de pruebas:} 1 — \textbf{Estado:} 1/1 PASANDO (100\,\%).

\subsubsection*{Cómo se implementó}

\begin{lstlisting}[caption={\texttt{Pago.generarComprobanteDigitalID} (\texttt{src/core/domain/models/Pago.js})}]
generarComprobanteDigitalID(casa, numeroSecuencial) {
    const hoy = new Date();
    const anio = hoy.getFullYear();
    const mes = String(hoy.getMonth() + 1).padStart(2, '0');
    const casaSanitizada = casa.replace(/[^a-zA-Z0-9]/g, '').toUpperCase();
    const secuencia = String(numeroSecuencial).padStart(4, '0');
    this.reciboId = `AGE-${anio}-${mes}-${casaSanitizada}-${secuencia}`;
    return this.reciboId;
}
\end{lstlisting}

\subsubsection*{Resultados Obtenidos}
La prueba \emph{RF-3.2/RF-3.4: permite aprobar desde la API y genera recibo inmutable}
de \texttt{pagoController.test.js} aprueba un pago vía \texttt{POST /api/pagos/:id/aprobar}
y verifica que se genere el identificador de recibo con el formato
\texttt{AGE-YYYY-MM-CASA-SEQ} y que el pago quede en estado \texttt{APROBADO} e inmutable.
El formato también se corrobora dentro de la prueba de amortización FIFO de la
Sección~4.11. La verificación binaria byte a byte del PDF (\texttt{generarPdfRecibo}) queda
fuera del alcance automatizado por tratarse de comparación de contenido binario, tal como
se indica en la Sección~2.1.

% ----------------------------------------------------------------------------
\subsection{REQ013: Notificaciones Automáticas de Pago}
% ----------------------------------------------------------------------------
\textbf{Archivo:} \texttt{tests/core/business/services/NotificationService.test.js} (3 casos),
complementado por la verificación con \emph{spy} dentro de \texttt{PagoService.test.js}.\\
\textbf{Total de pruebas:} 3 — \textbf{Estado:} 3/3 PASANDO (100\,\%).

\subsubsection*{Cómo se implementó}

\begin{lstlisting}[caption={Cola de notificaciones con reintentos (\texttt{src/core/business/services/NotificationService.js})}]
async procesarCola() {
    if (this.procesando) return;
    this.procesando = true;

    while (this.cola.length > 0) {
        const item = this.cola[0];
        try {
            item.intentos++;
            await this.simularApiWhatsApp(item.telefono, item.mensaje);
            this.cola.shift();
        } catch (err) {
            if (item.intentos >= 3) {
                await pagoRepository.registrarAuditoria("NOTIFICACION_WHATSAPP_FALLO_PERMANENTE", { telefono: item.telefono, intentos: item.intentos });
                this.cola.shift();
            } else {
                const esTest = process.env.NODE_ENV === 'test';
                await new Promise(resolve => setTimeout(resolve, esTest ? 0 : 30000)); // RF-04: reintento cada 30s
            }
        }
    }
    this.procesando = false;
}
\end{lstlisting}

\begin{lstlisting}[caption={Verificación indirecta con \texttt{jest.spyOn} (\texttt{PagoService.test.js})}]
const spyNotif = jest.spyOn(notificationService, 'enviarNotificacionPago');
const resultado = await pagoService.validarAprobarPago(pagoId, fechaReferencia);
// ...
expect(spyNotif).toHaveBeenCalled();
spyNotif.mockRestore();
\end{lstlisting}

\subsubsection*{Resultados Obtenidos}
El archivo dedicado \texttt{NotificationService.test.js} ejercita tres escenarios de forma
aislada: (i) registro de una notificación exitosa con tipo, destinatario, fecha e intentos;
(ii) advertencia persistente cuando el copropietario no tiene teléfono registrado; y
(iii) reintento de la cola hasta tres veces conservando el mensaje pendiente. Además, la
prueba con \texttt{jest.spyOn} en \texttt{PagoService.test.js} confirma que
\texttt{validarAprobarPago} invoca al servicio de notificaciones tras aprobar un pago. La
entrega real por WhatsApp Business API queda fuera del alcance por requerir credenciales de
despliegue.

% ----------------------------------------------------------------------------
\subsection{Pruebas Transversales de Modelo de Dominio y Objetos de Valor}
% ----------------------------------------------------------------------------
\textbf{Archivos:} \texttt{Usuario.test.js} (7), \texttt{Cedula.test.js} (7) y
\texttt{Telefono.test.js} (7).\\
\textbf{Total de pruebas:} 21 — \textbf{Estado:} 21/21 PASANDO (100\,\%).

Estos tres componentes no pertenecen a un único requisito, sino que constituyen la
infraestructura de validación reutilizada por \texttt{REQ001} (formato de usuario y
contraseña), \texttt{REQ004} (validación de cédula/teléfono al importar) y
\texttt{REQ007} (validación al modificar copropietarios).

\subsubsection*{Cómo se implementó}

\begin{lstlisting}[caption={Algoritmo Módulo 10 ecuatoriano (\texttt{src/core/domain/valueObjects/Cedula.js})}]
validarCedulaEcuatoriana(cedula) {
    if (cedula.length !== 10) return false;
    const digitoRegion = parseInt(cedula.substring(0, 2), 10);
    if ((digitoRegion < 1 || digitoRegion > 24) && digitoRegion !== 30) return false;

    const ultimoDigito = parseInt(cedula.substring(9, 10), 10);
    let suma = 0;
    for (let i = 0; i < 9; i++) {
        let multiplicador = (i % 2 === 0) ? 2 : 1;
        let componente = parseInt(cedula.substring(i, i + 1), 10) * multiplicador;
        if (componente > 9) componente -= 9;
        suma += componente;
    }
    const digitoVerificador = (suma % 10 === 0) ? 0 : 10 - (suma % 10);
    return digitoVerificador === ultimoDigito;
}
\end{lstlisting}

Para probar el caso de provincia ``consular exterior'' (30), fue necesario construir
manualmente una cédula válida que superara el algoritmo, documentando el cálculo
directamente en el comentario de la prueba (ver también Sección~6.6):

\begin{lstlisting}[caption={Cédula exterior válida calculada manualmente (\texttt{Cedula.test.js})}]
test('Debería permitir provincia exterior 30', () => {
    // Dígitos: 3 0 0 0 0 0 0 0 1 2 -> Posiciones impares x2: 3*2=6, 0, 0, 0, 1*2=2 -> Suma = 8
    // Posiciones pares x1: 0, 0, 0, 0 -> Suma = 0. Total = 8. Verificador: 10 - 8 = 2.
    // 3000000012 es una cédula exterior válida.
    const ced = new Cedula('3000000012');
    expect(ced.valor).toBe('3000000012');
});
\end{lstlisting}

\subsubsection*{Resultados Obtenidos}
\begin{itemize}[nosep]
    \item \texttt{Cedula}: validación del algoritmo Módulo 10, limpieza de guiones/espacios,
    rechazo de provincia inválida y aceptación de la provincia exterior 30 — verificado.
    \item \texttt{Telefono}: normalización a formato internacional
    (\texttt{+593}\ldots) desde los cuatro formatos de entrada aceptados (local con/sin
    cero, internacional con/sin signo \texttt{+}) y rechazo de números fijos — verificado.
    \item \texttt{Usuario}: longitud de usuario (3–20 caracteres), formato de correo,
    política de complejidad de contraseña (mayúscula, minúscula, número, carácter especial)
    y excepción para contraseñas ya cifradas con \texttt{bcrypt} (prefijo
    \texttt{\$2b\$}) — verificado.
    \item \texttt{Usuario.esBloqueado()}: estado activo, bloqueo indefinido y bloqueo
    temporal vigente/expirado — verificado.
\end{itemize}

\clearpage

% ============================================================================
\section{Resultados Obtenidos}
% ============================================================================

\subsection{Resumen de Ejecución}

\textbf{Fecha de ejecución:} 23 de julio de 2026.\quad
\textbf{Ambiente:} Windows 11, Node.js, \texttt{jest --runInBand}.\quad
\textbf{Resultado:} TODOS LOS TESTS PASANDO.

\begin{table}[H]
\centering
\caption{Resumen global de ejecución de la suite de pruebas}
\label{tab:resumen-ejecucion}
\begin{tabular}{lc}
\toprule
\textbf{Métrica} & \textbf{Valor} \\
\midrule
Suites de prueba & 17 / 17 (100\,\%) \\
Casos de prueba & 86 / 86 (100\,\%) \\
Tiempo de ejecución & \textasciitilde 5\,s (posterior al primer arranque) \\
Requisitos con cobertura completa & 12 / 13 \\
Requisitos con cobertura parcial & 1 / 13 (REQ006) \\
Snapshots & 0 \\
\bottomrule
\end{tabular}
\end{table}

\subsection{Desglose por Archivo}

\begin{table}[H]
\centering
\caption{Desglose de pruebas por archivo}
\label{tab:desglose-archivo}
\begin{tabular}{llcc}
\toprule
\textbf{Categoría} & \textbf{Archivo} & \textbf{Tests} & \textbf{Estado} \\
\midrule
Patrones de Diseño & \texttt{auditDecorator.test.js} & 2 & PASS \\
Patrones de Diseño & \texttt{excelAdapter.test.js} & 4 & PASS \\
Patrones de Diseño & \texttt{moraStrategy.test.js} & 4 & PASS \\
Servicios de Negocio & \texttt{AuthService.test.js} & 9 & PASS \\
Servicios de Negocio & \texttt{CopropietarioService.test.js} & 7 & PASS \\
Servicios de Negocio & \texttt{ExcelImportSample.test.js} & 1 & PASS \\
Servicios de Negocio & \texttt{NotificationService.test.js} & 3 & PASS \\
Servicios de Negocio & \texttt{PagoService.test.js} & 10 & PASS \\
Modelo de Dominio & \texttt{Usuario.test.js} & 7 & PASS \\
Objetos de Valor & \texttt{Cedula.test.js} & 7 & PASS \\
Objetos de Valor & \texttt{Telefono.test.js} & 7 & PASS \\
Datos (repositorio) & \texttt{sessionRepository.test.js} & 3 & PASS \\
Presentación (integración) & \texttt{authController.test.js} & 4 & PASS \\
Presentación (integración) & \texttt{copropietarioController.test.js} & 5 & PASS \\
Presentación (integración) & \texttt{pagoController.test.js} & 4 & PASS \\
Presentación (\emph{middleware}) & \texttt{authMiddleware.test.js} & 6 & PASS \\
Presentación (vista) & \texttt{dashboard.test.js} & 3 & PASS \\
\midrule
\textbf{TOTAL} & & \textbf{86} & \textbf{PASS} \\
\bottomrule
\end{tabular}
\end{table}

\subsection{Métricas por Requisito}

\begin{table}[H]
\centering
\caption{Métricas de cobertura por requisito funcional}
\label{tab:metricas-req}
\begin{tabular}{lp{6.0cm}cl}
\toprule
\textbf{REQ} & \textbf{Nombre} & \textbf{Tests} & \textbf{Cobertura} \\
\midrule
REQ001 & Gestión de Acceso y Registro de Usuarios & 12 & Completa \\
REQ002 & Recuperación de Contraseña & 4 & Completa \\
REQ003 & Configuración de Perfiles & 6 & Completa \\
REQ004 & Importación de Copropietarios & 8 & Completa \\
REQ005 & Consulta de Copropietarios & 3 & Completa \\
REQ006 & Generación de Reportes & 0\textsuperscript{*} & Parcial \\
REQ007 & Modificación de Copropietarios & 3 & Completa \\
REQ008 & Eliminación de Copropietarios & 3 & Completa \\
REQ009 & Registro de Pago de Expensa & 4 & Completa \\
REQ010 & Validación de Pagos & 5 & Completa \\
REQ011 & Cálculo de Sanción por Mora & 10 & Completa \\
REQ012 & Generación de Comprobante de Pago & 1 & Completa \\
REQ013 & Notificaciones Automáticas de Pago & 3 & Completa \\
Transversal & Modelo de Dominio / Objetos de Valor & 21 & Completa \\
Regresión & Vista \texttt{dashboard.html} (enlaces/ObjectId) & 3 & Completa \\
\midrule
\textbf{TOTAL} & & \textbf{86} & \\
\bottomrule
\end{tabular}
\\[0.2cm]
{\footnotesize \textsuperscript{*} La restricción de reportes por rol se ejercita de forma
indirecta a través de las pruebas de integración de controladores; \texttt{obtenerReportePagos}
y \texttt{generarPdfReportePagos} no cuentan aún con un caso de prueba dedicado
(ver Sección~7.4).}
\end{table}

\subsection{Distribución de Pruebas por Requisito}

\begin{table}[H]
\centering
\caption{Distribución porcentual de pruebas por requisito}
\label{tab:distribucion}
\begin{tabular}{llrl}
\toprule
\textbf{REQ} & \textbf{Descripción breve} & \textbf{\%} & \textbf{Proporción} \\
\midrule
REQ001 & Acceso y Registro & 14.0\,\% & \rule{56pt}{6pt} \\
REQ002 & Recuperar Contraseña & 4.7\,\% & \rule{19pt}{6pt} \\
REQ003 & Perfiles & 7.0\,\% & \rule{28pt}{6pt} \\
REQ004 & Importación & 9.3\,\% & \rule{37pt}{6pt} \\
REQ005 & Consulta & 3.5\,\% & \rule{14pt}{6pt} \\
REQ007 & Modificación & 3.5\,\% & \rule{14pt}{6pt} \\
REQ008 & Eliminación & 3.5\,\% & \rule{14pt}{6pt} \\
REQ009 & Registro de Pago & 4.7\,\% & \rule{19pt}{6pt} \\
REQ010 & Validación de Pagos & 5.8\,\% & \rule{23pt}{6pt} \\
REQ011 & Mora / Expensas & 11.6\,\% & \rule{46pt}{6pt} \\
REQ012 & Comprobante de Pago & 1.2\,\% & \rule{5pt}{6pt} \\
REQ013 & Notificaciones & 3.5\,\% & \rule{14pt}{6pt} \\
Transversal & Dominio / Valor & 24.4\,\% & \rule{98pt}{6pt} \\
Regresión & Vista dashboard & 3.5\,\% & \rule{14pt}{6pt} \\
\bottomrule
\end{tabular}
\end{table}

\clearpage

% ============================================================================
\section{Problemas y Soluciones}
% ============================================================================

\subsection{Recalculo del Recargo por Mora}

\textbf{Problema:} la implementación original de \texttt{EcuandorianMora12Strategy}
calculaba un interés \emph{prorrateado diariamente} sobre una tasa anual del 12\,\%
(\texttt{monto * (0.12/365) * diasRetraso}), lo cual no coincidía con la regla de negocio
especificada: un recargo \textbf{fijo} del 12\,\% aplicado una sola vez por período
vencido, independientemente de los días de atraso transcurridos.

\begin{lstlisting}[caption={Fórmula anterior (incorrecta respecto a la regla de negocio)}]
// Antes
calcularMora(monto, diasRetraso) {
    if (diasRetraso <= 0) return 0;
    const tasaDiaria = 0.12 / 365;
    return parseFloat((monto * tasaDiaria * diasRetraso).toFixed(2));
}
\end{lstlisting}

\textbf{Solución:} se sustituyó la fórmula por un recargo fijo y se recalcularon a mano
los valores esperados en \texttt{moraStrategy.test.js} y en la prueba de amortización FIFO
de \texttt{PagoService.test.js} (de \texttt{\$100.33} de deuda con mora a
\texttt{\$112.00}, y del saldo pendiente resultante de \texttt{\$50.33} a
\texttt{\$62.00}), documentando la aritmética directamente en los comentarios de la
prueba para que quede trazable ante futuras revisiones.

\subsection{Ausencia de Generación de Deudas Mensuales}

\textbf{Problema:} el repositorio \texttt{pagoRepository.createDeuda} existía, pero
\textbf{ningún controlador ni servicio lo invocaba en producción}; solo se llamaba desde
los propios archivos de prueba para fabricar datos de prueba manualmente. Esto significaba
que, en un entorno real, jamás se generaba una deuda mensual real y la lógica de
amortización FIFO nunca tenía un caso genuino que resolver.

\textbf{Solución:} se implementó \texttt{PagoService.generarExpensasMensuales}, expuesto
mediante \texttt{POST /api/pagos/generar-expensas} (solo Administrador), y se añadieron
4 pruebas nuevas que verifican la generación, la no duplicación por período y el cálculo
correcto del vencimiento (incluido el cambio de año en diciembre).

\subsection{Redondeo de Punto Flotante en Cálculos Monetarios}

\textbf{Problema:} las operaciones aritméticas con montos en JavaScript (\texttt{100 *
0.12}, \texttt{150.75 * 0.12}) pueden introducir errores de precisión de punto flotante
(por ejemplo, \texttt{0.1 + 0.2 !== 0.3}), lo que provoca fallos de aserción
intermitentes si no se redondea de forma explícita.

\textbf{Solución:} todos los cálculos monetarios se normalizan con
\texttt{parseFloat(valor.toFixed(2))} antes de comparase, tanto en el código de
producción (\texttt{moraStrategy.js}, \texttt{PagoService.js}) como en las aserciones de
prueba, garantizando que las comparaciones sean deterministas a dos decimales.

\subsection{Determinismo Temporal en Pruebas de Mora}

\textbf{Problema:} \texttt{validarAprobarPago} calcula los días de atraso comparando la
fecha de vencimiento contra la fecha actual (\texttt{new Date()}); si la prueba dependiera
de la fecha real del sistema, el resultado (días de atraso, y por lo tanto la lista de
deudas vencidas) cambiaría cada vez que se ejecutara la suite en un día distinto,
produciendo pruebas \emph{no deterministas}.

\textbf{Solución:} el método acepta un parámetro opcional \texttt{fechaReferencia}, que
las pruebas fijan explícitamente (\texttt{new Date('2026-07-20')}) para congelar el
cálculo de días de atraso y obtener siempre el mismo resultado, independientemente del
día real en que se ejecute la suite.

\subsection{Verificación de un Efecto Secundario Asíncrono en Segundo Plano}

\textbf{Problema:} \texttt{NotificationService.enviarNotificacionPago} encola el mensaje y
dispara \texttt{procesarCola()} \textbf{sin esperar su resolución}
(\texttt{this.procesarCola();} sin \texttt{await}), de modo que el envío real ocurre en
segundo plano después de que la prueba ya continuó su ejecución. Un \texttt{await}
directo sobre el resultado del envío no reflejaría el comportamiento real del sistema.

\textbf{Solución:} en lugar de esperar el efecto secundario completo, se utilizó
\texttt{jest.spyOn(notificationService, 'enviarNotificacionPago')} para verificar
\'unicamente que el método de encolado fue invocado por \texttt{validarAprobarPago}, que es
la responsabilidad que le corresponde al servicio de pagos; la entrega real por WhatsApp
queda fuera del alcance de esa prueba (ver también la brecha de cobertura de REQ013 en
la Sección~4.13).

\subsection{Construcción de Datos de Prueba Válidos para el Algoritmo Módulo 10}

\textbf{Problema:} para probar el caso de la provincia ``consular exterior'' (código 30)
del algoritmo de validación de cédula ecuatoriana, se requería un número de 10 dígitos que
superara el dígito verificador Módulo 10; generar uno al azar y ajustarlo por prueba y
error hubiera sido poco confiable y difícil de mantener.

\textbf{Solución:} se calculó manualmente un valor válido (\texttt{3000000012})
aplicando el propio algoritmo paso a paso, y ese cálculo se documentó como comentario
dentro de la prueba (ver Sección~4.14), de modo que cualquier persona que revise el test
pueda verificar la aritmética sin tener que ejecutar el código.

\subsection{Tiempo de Arranque de \texttt{mongodb-memory-server}}

\textbf{Problema:} la primera ejecución de la suite tras una instalación limpia del
proyecto tarda considerablemente más (\textasciitilde 18\,s) que las ejecuciones
posteriores, porque \texttt{mongodb-memory-server} debe descargar y almacenar en caché el
binario de MongoDB correspondiente al sistema operativo.

\textbf{Solución:} no se requirió ningún cambio de código; se documenta el comportamiento
esperado para que no se interprete como una regresión de rendimiento. En ejecuciones
subsecuentes, con el binario ya cacheado localmente, el tiempo total baja a
\textasciitilde 5\,segundos.

\clearpage

% ============================================================================
\section{Conclusiones}
% ============================================================================

\subsection{Logros Alcanzados}
\begin{itemize}[nosep]
    \item 100\,\% de los casos de prueba existentes pasando (86/86) en 17 suites.
    \item 12 de 13 requisitos funcionales del backlog con cobertura automatizada completa;
    el restante (REQ006, generación de reportes) cuenta con cobertura parcial documentada y
    trazable.
    \item Corrección verificada de una regla de negocio crítica (recargo de mora fijo del
    12\,\%) mediante pruebas de regresión.
    \item Incorporación de una funcionalidad ausente (generación de expensas mensuales)
    junto con su suite de pruebas correspondiente.
    \item Suite de ejecución rápida y reproducible (\textasciitilde 5\,s en caliente).
\end{itemize}

\subsection{Calidad del Software}

\textbf{Confiabilidad (5/5):} 0 pruebas intermitentes observadas en múltiples corridas;
resultados reproducibles gracias al aislamiento por \texttt{beforeEach} y a fechas de
referencia fijas en los cálculos de mora.\\
\textbf{Cobertura de Requisitos (4/5):} 12 de 13 requisitos con cobertura completa; se
documenta explícitamente la única brecha restante (REQ006) en lugar de omitirla.\\
\textbf{Mantenibilidad (5/5):} pruebas organizadas por capa arquitectónica
(dominio/objetos de valor, servicios, patrones, presentación), con comentarios que
documentan la aritmética esperada en los casos financieros.\\
\textbf{Rendimiento (5/5):} suite completa en segundos, apta para ejecutarse en cada
cambio antes de integrar código nuevo.

\subsection{Lecciones Aprendidas}
\begin{enumerate}[nosep]
    \item \textbf{Las reglas de negocio financieras deben verificarse contra la
    especificación, no solo contra el código existente:} el error de la fórmula de mora
    (interés prorrateado vs.\ recargo fijo) habría pasado inadvertido si las pruebas solo
    hubiesen validado que el código hiciera lo que ya hacía.
    \item \textbf{Una funcionalidad puede estar ``implementada'' en el repositorio de datos
    sin estar realmente disponible en producción} si ningún controlador o servicio la
    invoca; las pruebas de integración end-to-end (aunque sea a nivel de servicio) ayudan a
    detectar estas rutas muertas.
    \item \textbf{El determinismo temporal es esencial en dominios financieros:} cualquier
    cálculo que dependa de \texttt{new Date()} debe aceptar una fecha de referencia
    inyectable para que las pruebas sean reproducibles indefinidamente.
    \item \textbf{No toda dependencia externa debe mockearse igual:} aislar por completo
    con \texttt{jest.mock()} es adecuado para middlewares de control de flujo, pero para
    lógica que depende de consultas relacionales reales, una base de datos en memoria
    produce pruebas más representativas del comportamiento en producción.
    \item \textbf{Documentar la aritmética esperada dentro del propio test} (como se hizo
    en los casos de mora y del algoritmo Módulo 10) reduce el costo de mantenimiento futuro
    de las pruebas.
\end{enumerate}

\subsection{Recomendaciones Futuras}
\begin{enumerate}[nosep]
    \item \textbf{Cerrar la brecha de cobertura restante:} añadir pruebas dedicadas para
    \texttt{obtenerReportePagos}/\texttt{generarPdfReportePagos} (REQ006), única línea del
    backlog con cobertura sólo indirecta. Como refuerzo, ampliar
    \texttt{NotificationService.test.js} para ejercitar el registro de auditoría por fallo
    permanente de la cola de WhatsApp.
    \item \textbf{Pruebas de integración HTTP adicionales:} extender
    \texttt{pagoController} con \texttt{supertest} para verificar la carga real de
    comprobantes vía \texttt{multer} (tipo de archivo, límite de 5\,MB) y el flujo completo
    de aprobación/rechazo por HTTP.
    \item \textbf{Cobertura de código automatizada:} incorporar \texttt{--coverage} en el
    script de \texttt{npm test} y fijar un umbral mínimo por carpeta (dominio, servicios,
    presentación) para detectar regresiones de cobertura en el tiempo.
    \item \textbf{Pruebas end-to-end de la interfaz visual:} usar Playwright o Cypress
    para validar el flujo completo desde \texttt{login.html} y \texttt{dashboard.html}
    hasta la API, incluyendo la expiración de sesión por inactividad.
    \item \textbf{Integración continua (CI/CD):} automatizar la ejecución de
    \texttt{npm test} en cada \emph{pull request} mediante GitHub Actions.
    \item \textbf{Pruebas de mutación:} evaluar herramientas como Stryker para medir la
    calidad real de las aserciones existentes, más allá del simple conteo de casos que
    pasan.
\end{enumerate}

\subsection{Resumen Final}

El plan de pruebas unitarias de AliGest cubre, con 86 casos automatizados distribuidos en
17 suites, el núcleo de la lógica de negocio del sistema: autenticación y control de
acceso, administración de copropietarios, y el ciclo completo de pagos, mora, auditoría y
notificaciones. Durante la elaboración de este informe se detectó y corrigió una
discrepancia real entre la fórmula de mora implementada y la regla de negocio especificada,
y se incorporó la funcionalidad de generación de expensas mensuales que faltaba en el flujo
de producción, ambas respaldadas por pruebas de regresión. La única brecha de cobertura
restante (la invocación directa de los reportes financieros, REQ006) se documenta de forma
explícita como trabajo pendiente, en lugar de presentarse como resuelta.

\textbf{Estado final: APROBADO CON OBSERVACIONES} — el núcleo funcional crítico
(autenticación, copropietarios, pagos, mora, comprobantes y notificaciones) está
completamente cubierto y verificado; sólo la invocación directa de los reportes financieros
(REQ006) requiere una prueba adicional antes de considerarse con cobertura total.

\clearpage

% ============================================================================
\section*{Anexos}
\addcontentsline{toc}{section}{Anexos}
% ============================================================================

\subsection*{Anexo A: Comandos de Ejecución}

\begin{lstlisting}[language=bash,caption={Comandos útiles para ejecutar la suite}]
# Ejecutar todos los tests (según package.json: jest --runInBand)
npm test

# Ejecutar un archivo de prueba específico
npx jest tests/core/business/services/PagoService.test.js

# Ejecutar con salida detallada por caso
npx jest --verbose

# Ejecutar un caso puntual por nombre (coincidencia de texto)
npx jest -t "amortización FIFO"
\end{lstlisting}

\subsection*{Anexo B: Estructura de Prueba Típica}

\begin{lstlisting}[caption={Patrón Arrange-Act-Assert utilizado en todo el proyecto}]
describe('Componente - Funcionalidad', () => {
    beforeEach(async () => {
        // Arrange: preparar datos (fixtures) y limpiar la base de datos en memoria
    });

    test('Debe hacer algo específico', async () => {
        // Arrange: datos particulares del caso
        const datos = { /* ... */ };

        // Act: ejecutar la funcionalidad bajo prueba
        const resultado = await servicio.metodo(datos);

        // Assert: verificar el resultado esperado
        expect(resultado).toBe(/* valor esperado */);
    });
});
\end{lstlisting}

\subsection*{Anexo C: Configuración Completa}

\begin{lstlisting}[caption={\texttt{jest.config.js}}]
module.exports = {
  testEnvironment: 'node',
  verbose: true,
  setupFilesAfterEnv: ['./tests/setup.js'],
  testTimeout: 30000
};
\end{lstlisting}

\begin{lstlisting}[caption={Fragmento relevante de \texttt{package.json}}]
{
  "name": "aligest-prototype",
  "version": "1.0.0",
  "scripts": {
    "start": "node app.js",
    "test": "jest --runInBand"
  },
  "dependencies": {
    "bcrypt": "^5.1.1",
    "express": "^4.19.2",
    "mongoose": "^9.7.4",
    "multer": "^2.2.0",
    "pdfkit": "^0.19.1",
    "xlsx": "^0.18.5"
  },
  "devDependencies": {
    "jest": "^30.4.2",
    "mongodb-memory-server": "^11.2.0",
    "supertest": "^7.2.2"
  }
}
\end{lstlisting}

\begin{lstlisting}[caption={Mecanismo de expiración de sesión por inactividad (RNF-01), \texttt{authMiddleware.js}}]
const SESION_MAX_INACTIVIDAD_MS = 2 * 60 * 60 * 1000; // 2 horas

function verificarSesion(req, res, next) {
    const role = req.headers['x-user-role'];
    const userId = req.headers['x-user-id'];
    if (!role) {
        return res.status(401).json({ error: "No autorizado. Inicie sesión para acceder a AliGest." });
    }
    if (role !== 'ADMINISTRADOR' && role !== 'COPROPIETARIO') {
        return res.status(403).json({ error: "Acceso denegado. Rol inválido o desconocido." });
    }

    const lastActivity = parseInt(req.headers['x-last-activity'], 10);
    if (!isNaN(lastActivity) && Date.now() - lastActivity > SESION_MAX_INACTIVIDAD_MS) {
        return res.status(401).json({
            error: "Sesión expirada por inactividad. Por favor, inicie sesión nuevamente.",
            sesionExpirada: true
        });
    }

    req.user = { role, id: userId || null };
    next();
}
\end{lstlisting}

\subsection*{Anexo D: Métricas Detalladas}

\begin{table}[H]
\centering
\caption{Distribución de pruebas por categoría arquitectónica}
\label{tab:metricas-categoria}
\begin{tabular}{lcc}
\toprule
\textbf{Categoría} & \textbf{Tests} & \textbf{\%} \\
\midrule
Modelo de dominio y objetos de valor & 21 & 24.4\,\% \\
Servicios de negocio & 30 & 34.9\,\% \\
Patrones de diseño (Strategy/Decorator/Adapter) & 10 & 11.6\,\% \\
Datos (repositorio de sesiones) & 3 & 3.5\,\% \\
Presentación (\emph{middlewares}, controladores y vistas) & 22 & 25.6\,\% \\
\midrule
\textbf{TOTAL} & \textbf{86} & \textbf{100\,\%} \\
\bottomrule
\end{tabular}
\end{table}

\begin{table}[H]
\centering
\caption{Distribución de pruebas por tipo de validación}
\label{tab:metricas-tipo}
\begin{tabular}{lc}
\toprule
\textbf{Tipo de validación} & \textbf{Tests aproximados} \\
\midrule
Creación y validación de objetos de dominio & 21 \\
Reglas de negocio y validación cruzada contra BD & 30 \\
Control de acceso y expiración de sesión & 13 \\
Patrón Decorator/Strategy/Adapter & 10 \\
Integración HTTP (\texttt{supertest}) & 9 \\
Regresión de vista y serialización de enlaces (\texttt{dashboard.html}) & 3 \\
\bottomrule
\end{tabular}
\end{table}

\bigskip
\noindent\rule{\textwidth}{0.4pt}

\begin{center}
\textbf{FIN DEL INFORME}\\[0.2cm]
\textit{Documento generado el 23 de julio de 2026}\\
\textit{Sistema AliGest — Plan de Pruebas Unitarias}\\
\textit{Requisitos REQ001–REQ005, REQ007–REQ013 y los bloques transversal (dominio/objetos
de valor) y de regresión de vista: cobertura completa (86 pruebas en 17 suites). Sólo
REQ006 mantiene cobertura parcial documentada.}
\end{center}

\end{document}
